% SGA 3, Expose V, Loop 1 English reconstruction.
% Scope: continuation of Theorem 4.1(iv), including Lemma 4.2 and its
% application; stops immediately before the commutative diagram on local p. 17.
% Authority: Polo--Gille Exp5-13oct24.pdf, component-local p. 16
% (combined-reader p. 278), re-edited page header 266; internal marker 268.
% Jacob Reinhold's CC BY 4.0 Markdown is comparison/drafting material only.

\phantomsection\label{proof:sga3-V-4.1-iv-affine}%
现在证明 \hyperref[thm:sga3-V-4.1-iv]{(iv)}. 由假设,
\[
  X_0=\operatorname{Spec}A_0,\qquad
  X_1=\operatorname{Spec}A_1,
\]
并且对 \(i=0,1\), 同态
\(\delta_i:A_0\to A_1\) 使 \(A_1\) 成为有限 \(A_0\)-模.
因此同态
\[
  A_0\otimes_BA_0\longrightarrow A_1
\]
更是有限的.

再假设
\[
  d=d_0\boxtimes d_1:X_1\longrightarrow X_0\times_YX_0
\]
是单态射. 由引理~\hyperref[lem:sga3-V-4.1.3]{4.1.3},
相应的同态 \(A_0\otimes_BA_0\to A_1\) 是满射.
我们将证明它是同构; 在证明过程中, 还将证明
\(p:X_0\to Y\) 有限局部自由.

只须证明, 对 \(B\) 的每个素理想 \(\mathfrak p\),
以 \(\delta_{0\mathfrak p}\) 和 \(\delta_{1\mathfrak p}\) 为分量的同态
\[
  (A_0)_{\mathfrak p}\otimes_{B_{\mathfrak p}}(A_0)_{\mathfrak p}
  \longrightarrow(A_1)_{\mathfrak p},
\]
是双射. 因此可假设 \(B\) 是局部环.
于是\hyperref[proof:sga3-V-4.1-b]{步骤~b)} 表明 \((A_0)_{\mathfrak p}\)
是半局部环. 事实上, 若 \(\mathfrak m\) 是它的一个极大理想,
则其余极大理想形如 \(\delta_0^{-1}(\mathfrak n)\),
其中 \(\mathfrak n\) 遍历 \(A_1\) 中满足
\(\delta_1^{-1}(\mathfrak n)=\mathfrak m\) 的素理想.
这样的素理想 \(\mathfrak n\) 至多有
\[
  n=[A_1:A_0]
\]
个.

作一个忠实平坦基变换后, 还可假设 \(B\) 的剩余域是无限的.%
\footnote{\label{note:sga3-V-25}编者注: 见引理~\hyperref[lem:sga3-V-4.1.2]{4.1.2}.}
于是可用\hyperref[lem:sga3-V-4.2]{下面的引理}.

\medskip
\phantomsection\label{lem:sga3-V-4.2}%
\noindent\textbf{引理 4.2.}---
\textit{设 \(B\) 是剩余域无限的局部环, \(A\) 是半局部环,
\(i:B\to A\) 是将 \(B\) 的极大理想 \(\mathfrak n\)
映入 \(A\) 的根基 \(\mathfrak r\) 的同态.
设 \(M\) 是秩为 \(n\) 的自由 \(A\)-模,
\(N\) 是 \(M\) 的一个 \(B\)-子模, 它作为 \(A\)-模生成 \(M\).
则 \(N\) 包含 \(M\) 的一组 \(A\)-基.}

回忆, \(M\) 的元素列 \(m_1,\ldots,m_n\) 是 \(A\)-基,
当且仅当它们在 \(M/\mathfrak rM\) 中的典范像构成
\(A/\mathfrak r\) 上的一组基. 因此可将
\[
  M\text{ 改为 }M/\mathfrak rM,\quad
  N\text{ 改为 }N/(N\cap\mathfrak rM),\quad
  A\text{ 改为 }A/\mathfrak r,\quad
  B\text{ 改为 }B/\mathfrak n.
\]
在此情形下, \hyperref[lem:sga3-V-4.2]{引理}是初等的. 具体地说, 若
\[
  A=K_1\times\cdots\times K_r,
  \qquad
  M=K_1^n\times\cdots\times K_r^n,
\]
则对每个 \(j\), 选取一个第 \(j\) 个分量非零的元素 \(x_j\in N\).
由于域 \(B\) 无限, 各 \(x_j\) 的一个适当的 \(B\)-线性组合 \(x\)
的所有分量均非零. 将 \(M\) 改为 \(M/Ax\), 对 \(n\) 归纳即可.

将\hyperref[lem:sga3-V-4.2]{此引理}应用于 \(A=A_0\), \(i:B\hookrightarrow A_0\),
通过 \(\delta_1\) 看成 \(A_0\)-模的 \(M=A_1\),
以及 \(N=\delta_0(A_0)\). 事实上, 由于
\[
  d_0\boxtimes d_1:X_1\longrightarrow X_0\times_YX_0
\]
是闭浸入, 以 \(\delta_0,\delta_1\) 为分量的同态
\[
  A_0\otimes_BA_0\longrightarrow A_1
\]
是满射. 这恰好是说 \(\delta_0(A_0)\) 生成 \(A_0\)-模 \(A_1\).

选取 \(a_1,\ldots,a_n\in A_0\), 使得
\(\delta_0(a_1),\ldots,\delta_0(a_n)\) 构成 \(A_1\)
的一组 \(A_0\)-基. 若 \(a_1,\ldots,a_n\) 构成 \(A_0\)
的一组 \(B\)-基, 则 \(A_0\otimes_BA_0\to A_1\) 将基
\[
  (1\otimes a_i)_{1\leq i\leq n}
\]
映到基
\[
  (\delta_0(a_i))_{1\leq i\leq n},
\]
因而为双射. 因此, 若
\[
  \varepsilon:\mathbb Z^n\longrightarrow A_0
\]
是将 \(\mathbb Z^n\) 的标准基映到 \(a_1,\ldots,a_n\)
的交换群同态, 则只须证明以 \(i\) 和 \(\varepsilon\) 为分量的映射
\[
  B\otimes_{\mathbb Z}\mathbb Z^n\longrightarrow A_0
\]
是双射.

% Translation coverage: entire assigned component; adjoining components retain the original statement/proof seams.
% Source lines: 119; translated body lines before this audit comment: 114.
% Math preserved: 74 inline formulas, 13 display blocks, 0 equation environments; mathematical prose translated where present.
% Labels: 3/3; hyperrefs: 7/7; SGARefLinks: 0/0; inline targets: 0/0; footnotes: 1/1.
% Figure arguments preserved; all rendered captions and natural-language source locators translated.
% Unresolved TODOs: 0; compilation and visual inspection remain with the volume integrator.
